\documentclass[12pt,a4paper,twocolumn]{article}

% --- minimal packages for ICP ---
\usepackage[utf8]{inputenc}
\usepackage[T1]{fontenc}
\usepackage[id]{babel}
\usepackage{amsmath,amssymb}
\usepackage{graphicx}
\usepackage{geometry}
\usepackage{hyperref}
\usepackage{cite}
\usepackage{booktabs}
\usepackage{enumitem}
\usepackage{caption}

\geometry{top=2cm,left=2cm,right=2cm,bottom=2cm}
\hypersetup{colorlinks=true,linkcolor=black,urlcolor=blue,citecolor=blue}
\setlength{\columnsep}{0.5cm}

\title{\textbf{Deteksi Dini Deforestasi pada Hutan Lindung \\
Menggunakan U-Net dan Citra Sentinel-2} \\
\large Idea Concept Paper --- Deforest.id}
\author{Peneliti -- Deforest.id}
\date{}

\begin{document}
\maketitle
\thispagestyle{empty}

\section{Pendahuluan}
Indonesia memiliki hutan tropis terluas ketiga di dunia dengan
$\sim94$ juta hektar, namun kehilangan $\sim300$ ribu hektar hutan
alam per tahun \cite{fao2022sofo}. Kawasan hutan lindung -- yang
berfungsi sebagai sistem penyangga kehidupan (tata air, erosi,
kesuburan tanah) -- sangat rentan karena deforestasi skala kecil
($<$2~ha) secara kumulatif menyumbang $>$56\% total emisi karbon
hutan tropis \cite{xu2026nature}.

Masalah utama: sistem pemantauan eksisting (Global Forest Watch,
SIMONTANA) hanya mendeteksi pada resolusi 30~m atau lebih, sehingga
petak deforestasi kecil yang ekologis signifikan sering lolos.
Indonesia menggunakan ambang minimum 6,25~ha untuk pelaporan REDD+
\cite{unfccc2022indonesia}, jauh di atas ambang dampak ekologis
2~ha \cite{xu2026nature}. Dibutuhkan sistem yang mampu mendeteksi
deforestasi pada skala 0,5--2~ha secara otomatis dan real-time.

\section{Ide dan Kebaruan}
Kami mengusulkan pipeline deteksi deforestasi otomatis pada citra
Sentinel-2 (resolusi 10~m) menggunakan arsitektur U-Net, yang
dilengkapi kerangka ambang batas peringatan dini tiga level:
\textit{Waspada} ($\ge$0,64~ha), \textit{Siaga} ($\ge$2~ha),
\textit{Kritis} ($\ge$6,25~ha) dengan amplifikasi peringatan untuk
hutan lindung. Kebaruan terletak pada:
\begin{enumerate}[nosep]
    \item Integrasi pseudo-labeling NDVI temporal + U-Net tanpa
          anotasi manual
    \item Heuristik deteksi awan RGB-NDVI sebagai solusi atas QA60
          yang tidak berfungsi
    \item Kerangka ambang batas yang mensintesis tiga pilar:
          definisi FAO, temuan ekologis \cite{xu2026nature,
          taubert2018fragmentation}, dan regulasi Indonesia
\end{enumerate}

\section{Metodologi}
Pipeline terdiri dari enam tahap: (1) akuisisi 7 scene Sentinel-2
dari GEE (5 kanal: RGB+NIR+NDVI); (2) tiling $512\times512$ dan
chipping $64\times64$ dengan overlap 50\%; (3) deteksi awan dengan
heuristik RGB ($R{>}180\land G{>}180\land B{>}180\land
\text{NDVI}{<}0,10$); (4) pseudo-label deforestasi via delta NDVI
($\Delta\text{NDVI}<-0,15$) + morfologi; (5) pelatihan U-Net 5
kanal input, Dice Loss, 50 epoch; (6) inferensi dan sistem
peringatan. Setiap piksel $10\times10$ m = 0,01~ha, sehingga
konversi luas deforestasi: $A(\text{ha}) = N_{\text{piksel}} \times 0,01$.

\section{Hasil Awal dan Roadmap}
Model diuji pada 5.695 sampel dari 2 scene dengan IoU 0,7256,
Dice 0,8410, Precision 0,8340, Recall 0,8480. Model mampu
mendeteksi deforestasi pada skala 0,64~ha per chip, jauh di bawah
ambang nasional 6,25~ha. Roadmap pengembangan:

\begin{enumerate}[nosep]
    \item Integrasi data Sentinel-1 SAR untuk menembus tutupan awan
    \item Validasi lapangan pada 3 lokasi hutan lindung di Indonesia
    \item Deployment berbasis web dengan streaming GEE + inference
          on-the-fly
    \item Integrasi peta kawasan hutan lindung (KHLK) untuk
          amplifikasi peringatan otomatis
\end{enumerate}

\small
\bibliographystyle{ieeetr}
\bibliography{icp-refs}

\end{document}
